% resumo em português (150-500 palavras, parágrafo único, 3-5 palavras-chave)
% O texto abaixo está comentado; remova o "%" de cada linha para publicar.
\setlength{\absparsep}{18pt}
\begin{resumo}
% Este trabalho desenvolveu e avaliou um soft-sensor embarcado para estimar a
% potência ativa CA de inversores fotovoltaicos a partir de variáveis
% meteorológicas, cobrindo falhas de telemetria. Os dados vieram da Rede
% Solarimétrica de Pernambuco, com três inversores de 2~kW. O conjunto foi
% preparado em quatro granularidades (5~min, 10~min, 1~h e 1~d), com
% separação temporal entre treino, calibração e teste. Dois modelos foram
% comparados: Random Forest e perceptron de múltiplas camadas. Os modelos
% foram treinados, avaliados no teste de 2025 e embarcados em um ESP32-S3,
% com bibliotecas geradas por emlearn e TensorFlow Lite. No teste, o MLP
% apresentou MAE entre $0{,}160$ e $0{,}037$~kW e R$^2$ entre $0{,}808$ e
% $0{,}981$; o Random Forest ficou próximo, com MAE entre $0{,}164$ e
% $0{,}040$~kW e R$^2$ entre $0{,}805$ e $0{,}972$, e menor custo de
% memória. A análise cruzada mostrou que um modelo treinado em 1~h manteve
% R$^2$ de $0{,}784$ (MLP) e $0{,}777$ (RF) quando avaliado em 5~min. No
% ensaio na placa, o Random Forest de 1~h reproduziu o resultado
% \textit{offline}, com MAE de $0{,}066$~kW fora dos \textit{gaps} e
% $0{,}065$~kW dentro deles, contra $0{,}312$~kW da interpolação linear e
% $0{,}376$~kW da persistência. O resultado indica que o soft-sensor cobre a
% lacuna de telemetria com erro próximo do normal e muito abaixo das
% baselines, e que o caminho embarcado não degrada a inferência.
%
% \textbf{Palavras-chave:} aprendizado de máquina; sistemas fotovoltaicos;
% soft-sensor; ESP32-S3; dados faltantes.
\end{resumo}

% resumo em inglês
\begin{resumo}[Abstract]
\begin{otherlanguage*}{english}
% This work developed and evaluated an embedded soft sensor that estimates
% the AC active power of photovoltaic inverters from meteorological
% variables, covering telemetry failures. The data came from the
% Pernambuco Solarimetric Network, with three 2~kW inverters. The dataset
% was prepared at four granularities (5~min, 10~min, 1~h and 1~d), with a
% temporal split between training, calibration and test. Two models were
% compared: Random Forest and a multilayer perceptron. The models were
% trained, evaluated on the 2025 test set and deployed on an ESP32-S3, with
% libraries generated by emlearn and TensorFlow Lite. On the test set, the MLP
% reached an MAE between $0.160$ and $0.037$~kW and an R$^2$ between
% $0.808$ and $0.981$; the Random Forest was close, with an MAE between
% $0.164$ and $0.040$~kW and an R$^2$ between $0.805$ and $0.972$, and a
% lower embedded cost. The cross-granularity analysis showed that a model
% trained at 1~h kept an R$^2$ of $0.784$ (MLP) and $0.777$ (RF) when
% evaluated at 5~min. In the board test, the 1~h Random Forest reproduced
% the offline result, with an MAE of $0.066$~kW outside the gaps and
% $0.065$~kW inside them, against $0.312$~kW for linear interpolation and
% $0.376$~kW for persistence. The result indicates that the soft sensor
% covers the telemetry gap with an error close to the normal one and well
% below the baselines, and that the embedded path does not degrade the
% inference.
%
% \textbf{Keywords:} machine learning; photovoltaic systems; soft sensor;
% ESP32-S3; missing data.
\end{otherlanguage*}
\end{resumo}
